\phantomsection
\chapter*{Rahasia Sel Nomor 7}
\addcontentsline{toc}{section}{Rahasia Sel Nomor 7 --- Alban Ramdalova}
\penulis{Alban Ramdalova}

\awalcerpen{M}{alam di Lapas Sukamiskin} selalu berbau besi karatan dan tembok basah. Dion, cowok apes yang dipenjara gara-gara difitnah mencuri perhiasan pejabat korup, cuma bisa meringkuk di Sel Nomor 7. Sel itu kecil, pengap, dan kata napi tua adalah tempat pembuangan sial.


Di sini bukan cuma orang jahat yang harus kamu takutin, kata Pak Jaka, pria paruh baya yang sudah belasan tahun menghuni sel itu. ``Kalau lewat tengah malam ada yang mengetuk tembok tiga kali, jangan pernah dijawab. Kalau terdengar suara napas di atas kepala, lihat lantai saja, lalu berdoa.''


Di luar sel mereka, ada dua geng besar yang saling membenci. Geng Rajawali pimpinan Bang Tito yang bertubuh besar dengan tato naga di punggung, dan Geng Tengkorak punya Jarot si pria mata satu. Biasanya mereka hobi memeras jatah makan napi baru. Tapi malam itu, suasana mendadak berubah drastis.


Tepat jam dua subuh, udara sel seketika kerasa dingin sampai napas berembun. Lampu lorong kedap-kedip redup, lalu terdengar suara kuku menggesek jeruji besi.


Srak$\dots$ srak$\dots$


Dari ventilasi atas sel, turun sosok tinggi kurus dengan kulit sepucat mayat dan mata merah menyala. Itu sosok gaib yang paling ditakuti, Penghuni Sel Tua.


Bukannya melawan, dua bos geng yang biasanya garang malah panik setengah mati. Bang Tito yang berbadan bongsor langsung melompat spontan ke atas pangkuan anak buahnya sambil menjerit melengking. Sementara Jarot malah sibuk bersembunyi di balik tumpukan kasur kapuk sambil memeluk jimat. Sipir yang bertugas sengaja mengunci pintu lorong utama lalu kabur duluan.


Jangan pegang-pegang tangan gue, Tito! teriak Jarot dengan suara bergetar.


Gue takut, Rot! Itu setan mukanya mirip mantan bini gue kalau lagi ngamuk! Kata Tito dengan suara bergetar.


Dion yang sudah gemas melihat kelakuan konyol para bos geng itu akhirnya berteriak, ``Woy, Kalau kita cuma jerit-jeritan, kita mati semua malam ini, Ayo gabung!''


Akhirnya dua geng musuh itu terpaksa bekerja sama demi bertahan hidup. Aksi mereka dipenuhi komedi konyol. Bang Tito dijadikan benteng manusia di depan, Jarot sibuk melempar piring seng dan sapu ijuk ke arah hantu, sementara anak buahnya mengacungkan sendok makan seperti senjata sihir.


Di tengah kekacauan, Dion memperhatikan tingkah sang hantu. Tiap kali melayang mendekat, tangannya selalu menunjuk ke satu titik: lantai semen Sel Nomor 7 yang retak-retak.


Dia mau kasih tahu sesuatu! kata Dion.


Menggunakan sendok makan, potongan besi, dan mangkuk seng, kedua geng yang biasanya saling tebas kini bergotong-royong membongkar lantai semen tersebut. Di bawah tanah yang tergali, mereka menemukan sebuah kotak besi tua berlapis karat.


Saat kotak dibuka, di dalamnya tersimpan buku harian dan berkas lama. Terungkaplah kenyataan pahit: hantu itu dulunya adalah tahanan tak bersalah yang dibunuh dan mayatnya disembunyikan di balik tembok belasan tahun lalu. Ia mengamuk bukan karena jahat, melainkan hanya menuntut kebenaran atas kematiannya.


Dion segera mengambil radio komunikasi sipir yang tertinggal, lalu membacakan seluruh isi dokumen itu lewat radio penjara.


Tak lama setelah nama tahanan itu dibersihkan, wujud hantu yang tadinya menyeramkan perlahan memudar. Matanya tak lagi merah, berubah memancarkan cahaya halus, lalu menghilang tenang terbawa angin malam.


Pas pagi harinya ketika pintu penjara dibuka oleh sipir, pemandangannya bikin melongo. Bang Tito dan Jarot, dua musuh buyutan sedang tidur pulas saling tindih di sudut sel gara-gara kecapean. Sel Nomor 7 tak lagi angker, melainkan jadi saksi bisu bersatunya para napi demi sebuah keadilan.


\jedahias
